\claimheading{Induction Heads: General In-Context Learning}

\textbf{Source.} \emph{In-context Learning and Induction Heads} \citep{olsson2022context}.

\textbf{Description.} The two-head composition that performs in-context copying is advanced as the mechanistic source of in-context learning in general. In-context learning is operationalized as the difference in loss between the 500th and the 50th token of the context, and the claim rests on induction-head formation, the jump in that quantity, a loss bump and a PCA pivot all arriving together in a $2.5$--$5\times10^{9}$-token window. Evaluated on: the same 34 decoder-only models, with ablation confined to the twelve small models. Description mode: implementational-functional, carried to a computational-level capability by Argument 6.

\vspace{4pt}
{\footnotesize
\renewcommand{\arraystretch}{1.0}
\begin{longtable}{@{}llll>{\raggedright\arraybackslash}p{4.6cm}@{}}
\caption{\textbf{Induction heads: general in-context learning.} Full 36-criterion audit. Status: \textbf{C} = Confirmed,
\textbf{PC} = Partially confirmed, \textbf{U} = Untested, \textbf{I} = Inconclusive,
\textbf{D} = Disconfirmed, \textbf{N/A} = Not applicable.}\\
\toprule
\textbf{ID} & \textbf{Criterion} & \textbf{Status} & \textbf{Evidence} \\
\endfirsthead
\toprule
\textbf{ID} & \textbf{Criterion} & \textbf{Status} & \textbf{Evidence} \\
\endhead
\midrule
\multicolumn{4}{@{}l}{\textit{Construct Validity}} \\
C1  & Falsifiability            & C   & Argument 6 names the falsifier precisely, whatever the eventual outcome \\
C2  & Structural plausibility   & PC  & Weight-level for small attention-only models; the extension is analogy \\
C3  & Convergent validity       & PC  & Six lines of evidence, but at scale the paper's own table reads correlational \\
C4  & Discriminant validity     & D   & The title says general in-context learning; the measure is loss at token 500 \\
C5  & Nomological validity      & PC  & Connected in three directions, mostly by analogy rather than derivation \\
C6  & Complementation validity  & U   & The composition-head alternative is what a trans test would settle \\
\midrule
\multicolumn{4}{@{}l}{\textit{Measurement Validity}} \\
M1  & Reliability               & U   & One run per model, and at scale the evidence is timing across fifteen snapshots \\
M2  & Baseline separation       & PC  & The one-layer model is a clean negative for all four signature quantities \\
M3  & Stability                 & PC  & Free constants perturbed and the between-model picture survives \\
M4  & Calibration               & I   & Validated evaluators, but footnote 17 documents a sign-inverting artifact \\
M5  & Sensitivity               & U   & The only known-positive is literal copying in a two-layer model \\
M6  & Invariance                & PC  & The signature recurs across four series; a correlate's invariance is not the claim's \\
M7  & Selection correction      & U   & Top scorers examined on the data that ranked them, and the sample is biased \\
\midrule
\multicolumn{4}{@{}l}{\textit{Internal Validity}} \\
I1  & Necessity                 & I   & Ablation carries the metric in twelve models and the construct nowhere \\
I2  & Sufficiency               & U   & A semi-empirical argument stands in; nothing reconstructs the behavior \\
I3  & Minimality                & U   & Every ablation marginal and single-head, in the redundancy regime that defeats it \\
I4  & Specificity               & PC  & Ablating non-induction heads raises the score, which runs against the claim \\
I5  & Rival mechanism exclusion & D   & The composition rival is named twice by the authors and excluded neither time \\
I6  & Double dissociation       & U   & No second head population is defined that could be ablated against a second outcome \\
I7  & Confound control          & PC  & Exogenous confounds located outside the window; the shared latent is untouched \\
I8  & Confounding sensitivity   & U   & Every caution is verbal; the leading alternative is left unquantified \\
I9  & Epistatic interaction     & PC  & Composition is the subject and no joint ablation measures it \\
I10 & Rescue reversibility      & U   & The clean patterns are already cached; the restore costs one pass and is not run \\
I11 & Onset coupling            & C   & Four quantities move together in one token window, and an intervention shifts them \\
I12 & Offset coupling           & D   & The manipulation runs one way only; nothing suppresses formation and re-tests \\
\midrule
\multicolumn{4}{@{}l}{\textit{External Validity}} \\
E1  & Intervention reach        & PC  & Two levers of different kinds, neither reaching the models the claim is about \\
E2  & Prompt generalization     & C   & The defining stimulus is off-distribution by construction \\
E3  & Cross-task generalization & PC  & Three heads in one model on three tasks, labeled Plausibility by the authors \\
E4  & Cross-model generalization& D   & The signature holds across 34 models; nothing distinguishes the heads at scale \\
E5  & Graded response           & PC  & Continuous scores throughout and an all-or-nothing intervention \\
E6  & Novel prediction          & C   & The smeared-key architecture predicted the shift in advance, and it moved \\
\midrule
\multicolumn{4}{@{}l}{\textit{Interpretive Validity}} \\
V1  & Level declaration         & C   & Every sub-claim graded in a table, by reach and by model class \\
V2  & Level-evidence match      & D   & The assertion is mechanistic and unbounded where the tables say correlational \\
V3  & Alternative level         & PC  & The shared-latent alternative is stated precisely twice and tested never \\
V4  & Unlicensed labeling       & PC  & Both loaded terms hedged, and neither contrastively tested \\
V5  & Scope declaration         & C   & Scoped structurally, with the missing evidence graded rather than hedged \\
\midrule
\multicolumn{4}{@{}l}{\textbf{Total:} 6 Confirmed, 14 Partially confirmed, 2 Inconclusive, 9 Untested, 5 Disconfirmed} \\
\multicolumn{4}{@{}l}{\textbf{Verdict: Disconfirmed}} \\
\bottomrule
\end{longtable}}

\clearpage
\begin{table}[H]
\centering
\footnotesize
\caption{\textbf{Induction heads: general in-context learning.} Readings of the claim, from the version the authors state to the version the
field cites.}
\label{tab:induction_broad-readings}
\renewcommand{\arraystretch}{1.18}
\begin{tabular}{@{}>{\raggedright\arraybackslash}p{3.3cm}l>{\raggedright\arraybackslash}p{4.4cm}@{}}
\toprule
\textbf{Reading} & \textbf{Verdict} & \textbf{Missing} \\
\midrule
\textbf{Primary.} Induction heads are the mechanistic source of general in-context learning in transformers of any size
  & Disconfirmed
  & Nothing --- tested and failed: no ablation runs above the twelve small models, so nothing at scale separates induction heads from whatever else forms alongside them (E4, I5), and the adopted measure does not separate general in-context learning from the few-shot accuracy the field reads it as (C4) \\
\addlinespace[2pt]
The phase change at which induction heads form is when in-context learning arrives
  & Underdetermined
  & A test that separates the rival the authors state in Argument 1 --- that the phase change is when layer composition becomes available, forming induction heads and other composition mechanisms together --- from the causal reading (I5, I7, V3) \\
\addlinespace[2pt]
The share of in-context learning induction heads account for is the majority
  & Insufficient
  & No admissible measurement of the share: the paper's summary table asserts it while the semi-empirical argument offered for it states no fraction (I2), the ablation is single-head and all-or-nothing (E5), and the headline metric is written down five times with four of them reversing its sign (M4) \\
\addlinespace[2pt]
The copying the ablations cover is the literal case of a general retrieval mechanism
  & Underdetermined
  & The bridge is the in-context nearest-neighbor analogy, argued rather than derived (C2), and the paper's own strength table labels its large-model rows correlational and Argument 6 analogy (V2) \\
\bottomrule
\end{tabular}
\end{table}

\begin{table}[H]
\centering
\footnotesize
\caption{\textbf{Induction heads: general in-context learning.} Verdict: \textbf{Disconfirmed}. A dash means the tier is reached; a criterion at
Inconclusive, Untested or Disconfirmed blocks promotion.}
\label{tab:induction_broad-verdict}
\renewcommand{\arraystretch}{1.18}
\begin{tabular}{@{}l>{\raggedright\arraybackslash}p{3.6cm}>{\raggedright\arraybackslash}p{4.4cm}@{}}
\toprule
\textbf{Tier} & \textbf{Requires} & \textbf{Missing} \\
\midrule
Proposed & C1--C2; one admissible measurement & --- \\
\addlinespace[2pt]
Causally Suggestive & Necessity (I1); baselines (M2) & \textbf{I1} inconclusive: necessity holds for the metric in the twelve small models, is called only suggestive once MLPs are present, and is untested above them \\
\addlinespace[2pt]
Mechanistically Supported & Sufficiency (I2); reach (E1); specificity (I4) & \textbf{I2} unattempted: nothing reconstructs in-context learning from induction heads, and the argument offered instead states no fraction \\
\addlinespace[2pt]
Triangulated & Convergence (C3); replication (E2/E4); dissociation (I6) & \textbf{E4} Disconfirmed: the causal evidence never leaves the twelve small models; \textbf{C3}: only two of the six lines of evidence are interventional and both stop there \\
\addlinespace[2pt]
Validated & Calibration (M1--M7); interpretive audit (V1--V5); reach (E2, E4) & \textbf{V2}, \textbf{E4} Disconfirmed; \textbf{M4} inconclusive; \textbf{M1}, \textbf{M5}, \textbf{M7} unattempted, and the head pool is ranked and then sampled towards its high scorers \\
\bottomrule
\end{tabular}
\end{table}
